% ===================================================================
%  TABEL Nr. 9 — De berg. — De badplaats. — Het woud. — De jacht.
%  Traduction 2026 d'apres texto/io, controlee sur texto/fr et
%  texto/es. Consigne : voir texto/nl/10-tabel-01.tex.
% ===================================================================

%%K t09-apar-1 apar
\VUsaut{4.0mm}
\VUtitre{15.0pt}{60}{TABEL Nr. 9}
\VUsaut{3.0mm}

%%K t09-tit-1 sub
\VUcentre{12.0pt}{0}{\textit{Eerste tafereel.}}
\VUsaut{3.0mm}

%%K t09-tit-2 sub
\VUpk{11.4pt}{40}{De berg. --- De badplaats.}
\VUsaut{3.0mm}

%%K t09-c1-01-1 p
1. --- Sinds acht dagen ben ik in Prats. Ik kan niet de hele bewondering uitdrukken die ik
voelde toen ik voor het wonderbare \VUgras{gebergte} \textsuperscript{(1)} van de Pyreneeën
stond waarvan de \VUgras{kam} \textsuperscript{(2)} tegen de blauwe hemel als een reusachtige
feston uitkomt.

%%K t09-c1-02-1 p
2. --- Ik stelde mij niet voor, dat de \VUgras{toppen} \textsuperscript{(3)} van de Pyreneeën
zo'n grote \VUgras{hoogte} bereiken, en ik stond verstomd voor de prachtige \VUgras{gletsjers}
\textsuperscript{(5)} en de besneeuwde
\VUgras{bergpieken} \textsuperscript{(4)} waarop zulke schrikwekkende
\VUgras{lawines} rollen.

%%K t09-tit-3 sub
\VUsaut{3.0mm}
\VUpk{10.2pt}{40}{Berglandschap.}
\VUsaut{3.0mm}

%%K t09-c1-03-1 p
3. --- Al de dag na mijn aankomst ging ik naar de oude uitgedoofde
\VUgras{vulkaan} \textsuperscript{(6)} die bij Prats ligt. Op de dorre en kale
randen van de \VUgras{krater} ziet men nog goed de sporen van de doortocht van de
\VUgras{lava} die, verscheidene eeuwen geleden, langs de \VUgras{berghelling}
\textsuperscript{(7)} stroomde. Het lukte mij, op een van de \VUgras{hellingen}, enkele
brokstukken van die lava te vinden.

%%K t09-c1-04-1 p
4. --- Prats ligt aan de grens, en de \VUgras{vesting} \textsuperscript{(8)} die de
\VUgras{bergpas} \textsuperscript{(27)} verdedigt die Frankrijk van Spanje scheidt herinnert
helemaal aan die oude kastelen van de Rijnoever die van de top van hun rots een \VUgras{prooi}
schijnen te bespieden.

%%K t09-c1-05-1 p
5. --- Een enorme \VUgras{arend} \textsuperscript{(9)} die zijn nest niet ver van het
\VUgras{fort} \textsuperscript{(8)} heeft trekt dikwijls mijn aandacht. Die
\VUgras{roofvogel}, wiens \VUgras{snavel} sterk is, brengt zijn jongen dikwijls een lam of een
geitje dat hij met zijn \VUgras{klauwen} vasthoudt. Zo groot is de spanwijdte van zijn
\VUgras{vleugels} dat hij lang met zijn zware last kan zweven.

%%K t09-tit-4 sub
\VUsaut{3.0mm}
\VUpk{10.2pt}{40}{De bergwandelaar.}
\VUsaut{3.0mm}

%%K t09-c1-06-1 p
6. --- Daar is de aangenaamste wandeling de tocht naar de waterval \textsuperscript{(10)} van
« Cau ». De \VUgras{waterval} stort in kolken naar beneden en verdwijnt daarna als een mooi
\VUgras{beekje} in een \VUgras{dennenbos} \textsuperscript{(11)} ten westen van de stad. Wij
klommen door dat bos omhoog tot aan het \VUgras{weerkundig observatorium}
\textsuperscript{(12)} en van de top van het \VUgras{uitkijkpunt} zagen wij het hele
\VUgras{panorama} van de streek. Het \VUgras{landschap} is wonderbaar, en de \VUgras{natuur}
schijnt aan het land alle schatten kwistig te schenken waarover zij beschikt.

%%K t09-c1-07-1 p
7. --- Om af te dalen gebruikten wij de \VUgras{kabelspoorweg} \textsuperscript{(13)} en wij
hielden stil in een klein \VUgras{station} bij de
\VUgras{ruïnes} \textsuperscript{(14)} van een oud \VUgras{feodaal kasteel} dat
eertijds door de heren van Prats bewoond werd.

%%K t09-c1-08-1 p
8. --- Arm kasteel! Wat denkt het waarschijnlijk als het beneden de coquette
\VUgras{badplaats} \textsuperscript{(15)} ziet die nu een
\VUgras{modieuze kuuroord} is?

%%K t09-c1-08-2 p
Het \VUgras{klimaat} is zo aangenaam, het \VUgras{mineraalwater} zo werkzaam, dat de
\VUgras{kuur} die men in het \VUgras{sanatorium} \textsuperscript{(17)} toepast een waar
genoegen is.

%%K t09-tit-5 sub
\VUsaut{3.0mm}
\VUpk{10.2pt}{40}{Aangename kuurstad.}
\VUsaut{3.0mm}

%%K t09-c1-09-1 p
9. --- Overigens heeft men alles gedaan om het verblijf aangenaam te maken. Een groot
\VUgras{viaduct} \textsuperscript{(16)} werd gebouwd om een spoorweg mogelijk te maken; het
\VUgras{park} \textsuperscript{(18)} van de
\VUgras{kuurinrichting}, waar men \VUgras{concerten geeft}, is op bekoorlijke
wijze aangelegd. Verscheidene sierlijke « \VUgras{villa's} » \textsuperscript{(19)} werden om
het casino \textsuperscript{(21)} gebouwd, en een heel goed onderhouden \VUgras{straatweg}
\textsuperscript{(22)} maakt de verbindingen een stuk makkelijker.

%%K t09-c1-10-1 p
10. --- Een eenvoudig \VUgras{talud} \textsuperscript{(23)} scheidt de
\VUgras{weg} \textsuperscript{(20)} van het eigenlijke \VUgras{dal}
\textsuperscript{(25)}, in de bodem waarvan een sierlijk \VUgras{meertje}
\textsuperscript{(26)} ligt dat een helder \VUgras{beekje} \textsuperscript{(24)} voedt.

%%K t09-c1-11-1 p
11. --- Aan de andere kant van het dal wordt een \VUgras{ijzermijn} \textsuperscript{(28)}
ontgonnen. Verscheidene keren ging ik de \VUgras{mijnwerkers} zien die in de \VUgras{schacht}
afdaalden, en zelfs ging ik de \VUgras{gang} binnen waarin zij hun dag doorbrengen, terwijl
zij het \VUgras{erts} winnen dat het \VUgras{metaal} bevat.

%%K t09-c1-12-1 p
12. --- Een belangrijke \VUgras{ijzergieterij} werd bij de mijn gebouwd om de rijkdommen die
eruit gewonnen worden dadelijk te benutten. De
\VUgras{hoogoven} \textsuperscript{(29)}, verhit met de \VUgras{steenkool} die
hier overvloedig is, maakt het mogelijk snel en goedkoop \VUgras{gietijzer} te verkrijgen.

%%K t09-c1-13-1 p
13. --- Niet ver van de mijn graaft men een \VUgras{steengroeve} \textsuperscript{(30)} uit.
Noch \VUgras{graniet}, noch \VUgras{porfier}, noch zelfs \VUgras{zandsteen} hakt de oude
\VUgras{steenhouwer} met zijn
\VUgras{houweel}, maar eenvoudige \VUgras{bouwstenen} voor het bouwen van
huizen.

%%K t09-c1-14-1 p
14. --- Een van de mooiste sieraden van dat land zijn de \VUgras{dennen}
\textsuperscript{(31)}. Overal in de bodem van een \VUgras{ravijn} \textsuperscript{(32)}, aan
de oever van een \VUgras{bergstroom} \textsuperscript{(33)}, van een \VUgras{afgrond}
\textsuperscript{(34)} of van een steilte, geven hun dunne naalden een groene tint.

%%K t09-tit-6 sub
\VUsaut{3.0mm}
\VUpk{10.2pt}{40}{Te voet gaan.}
\VUsaut{3.0mm}

%%K t09-c1-15-1 p
15. --- Ik maak van mijn vakantie gebruik om \VUgras{lange} wandelingen te maken. De
\VUgras{wegen} \textsuperscript{(35)} die zich over de berg slingeren maken het mogelijk,
zonder op de afstand te letten, verscheidene \VUgras{kilometers} en dikwijls verscheidene
\VUgras{myriameters} af te leggen.

%%K t09-c1-15-2 p
\VUgras{Grenspalen} \textsuperscript{(36)} en \VUgras{wegwijzers}
\textsuperscript{(37)} tekenen de wegen waarop men dikwijls een jonge
\VUgras{bergbewoner} \textsuperscript{(38)} ontmoet met een \VUgras{knapzak}
\textsuperscript{(40)} aan zijn zij, die een \VUgras{marmot} \textsuperscript{(39)} draagt, of
een \VUgras{zigeuner} \textsuperscript{(41)} wiens \VUgras{bontmuts} \textsuperscript{(42)}
vreemd afsteekt bij de oude gescheurde \VUgras{kap} \textsuperscript{(43)}. Hij leidt aan een
\VUgras{ketting een beer} \textsuperscript{(44)} die afgericht is.

%%K t09-tit-7 sub
\VUpk{10.2pt}{40}{Bergbeklimming.}
\VUsaut{3.0mm}

%%K t09-c1-16-1 p
16. --- Bijna elke dag ga ik met een \VUgras{gids} \textsuperscript{(48)} een
\VUgras{beklimming} maken, en ik ben heel trots wanneer ik, met een \VUgras{zak}
\textsuperscript{(47)} op de rug, en de « \VUgras{alpenstok} » in de hand, als een echte
\VUgras{toerist} \textsuperscript{(46)} de \VUgras{rotsen} \textsuperscript{(45)} bestorm.

%%K t09-c1-17-1 p
17. --- Van de top van een berg lijken de bewoners van de vlakte niet dikker dan mieren; een
\VUgras{schaapskudde} \textsuperscript{(50)} die in een
\VUgras{weide} \textsuperscript{(49)} graast lijkt een kinderspeelgoed. Een
\VUgras{den} \textsuperscript{(51)} of ook een \VUgras{houten huisje}
\textsuperscript{(52)} lijkt nauwelijks een vlek op het groen.

%%K t09-c1-18-1 p
18. --- Tijdens die tochten ontmoeten wij dikwijls op het \VUgras{pad} \textsuperscript{(53)}
een \VUgras{gemsjager} \textsuperscript{(54)} die op zijn schouder het dier draagt dat hij
zojuist gedood heeft.

%%K t09-c1-19-1 p
19. --- Ik legde in mijn herbarium een heel merkwaardige tak \VUgras{varen}
\textsuperscript{(55)} en een mooi rozig takje \VUgras{heide} \textsuperscript{(57)} die ik
bij de ingang van een kleine \VUgras{grot} \textsuperscript{(56)} tijdens mijn laatste
wandeling plukte.

%%K t09-tit-8 sub
\VUsaut{3.0mm}
\VUcentre{12.0pt}{0}{\textit{Tweede tafereel.}}
\VUsaut{3.0mm}

%%K t09-tit-9 sub
\VUpk{11.4pt}{40}{Het woud. --- De jacht.}
\VUsaut{3.0mm}

%%K t09-c2-01-1 p
1. --- Gisteren bracht ik een heerlijke dag in het woud door. De bedrijvigheid die er heerst
onder de diertjes en de \VUgras{insecten} \textsuperscript{(5)} waardoor het bewoond wordt
boeide mij enorm.

%%K t09-c2-01-2 p
De \VUgras{spin} \textsuperscript{(1)} die haar \VUgras{web} \textsuperscript{(2)} tussen twee
takken weeft om de voorbijkomende
\VUgras{vlieg} \textsuperscript{(3)} te vangen, de \VUgras{slak}
\textsuperscript{(4)} die moeizaam haar huisje draagt, het vliegend hert dat op zoek is naar
zijn buit, de vlijtige mier die heel ijverig bij de
\VUgras{mierenhoop} \textsuperscript{(6)} bezig is, al die kleintjes gingen,
kwamen, werkten met buitengewone ijver.

%%K t09-c2-02-1 p
2. --- Een \VUgras{slang} \textsuperscript{(7)} die ik op het eerste gezicht voor een
\VUgras{adder} hield, maar die niets anders was dan een gewone
\VUgras{ringslang}, had zich onder een boomstam verscholen, en stak van tijd tot
tijd haar dunne kopje uit dat altijd klaar leek om te bijten. Voor mij kroop een dikke
\VUgras{naaktslak} \textsuperscript{(8)} zwaar voort nadat zij een van de smakelijke
\VUgras{aardbeitjes} \textsuperscript{(11)} verslonden had die daar overvloedig groeien.

%%K t09-c2-03-1 p
3. --- De grond was met \VUgras{bosbloemen} bekleed; \VUgras{sleutelbloemen}
\textsuperscript{(9)}, \VUgras{maagdenpalm} \textsuperscript{(10)}, en veel andere planten die
ik niet kende, bloeiden tussen \VUgras{graspollen} \textsuperscript{(12)}.

%%K t09-c2-04-1 p
4. --- In die streek is de \VUgras{truffel} bijna even gewoon als de
\VUgras{paddenstoel} \textsuperscript{(14)} en een \VUgras{big}
\textsuperscript{(13)} dat aan een boswachter toebehoorde, zocht gulzig of het er enkele zou
vinden.

%%K t09-tit-10 sub
\VUsaut{3.0mm}
\VUpk{10.2pt}{40}{Stroperij.}
\VUsaut{3.0mm}

%%K t09-c2-05-1 p
5. --- Ik woonde het \VUgras{einde} bij van een tafereel dat mij heel vermaakte. Met behulp
van de reuk van zijn \VUgras{dashond} \textsuperscript{(15)} had de
\VUgras{boswachter} \textsuperscript{(16)} zojuist een \VUgras{stroper}
\textsuperscript{(22)} verrast, op het ogenblik dat deze zich gereedmaakte om het
\VUgras{wild} \textsuperscript{(17)} mee te nemen dat hij gedood had. Omdat hij
liever zijn buit achterliet dan zich te laten aanhouden, was de stroper gevlucht, en een
\VUgras{haas} \textsuperscript{(18)}, een \VUgras{hinde} \textsuperscript{(19)}, een
\VUgras{ree} \textsuperscript{(20)} en een
\VUgras{everzwijn} \textsuperscript{(21)} lagen op het gras tot vreugde van de
boswachter.

%%K t09-c2-06-1 p
6. --- De verscheidenheid van de bomen die het woud bevat is verbazend. De
\VUgras{beuken} \textsuperscript{(23)} en de \VUgras{berken}
\textsuperscript{(26)}, de \VUgras{pijnbomen}, die de \VUgras{hars} \textsuperscript{(27)}
opleveren, de \VUgras{eiken} \textsuperscript{(29)} en ook veel andere vermengen hun takken.

%%K t09-c2-06-2 p
De \VUgras{klimop} \textsuperscript{(24)}, die zich aan de stam van de hoge bomen hecht,
kleurt het \VUgras{woud} \textsuperscript{(25)} met een glanzend groen, dat zich aangenaam
verenigt met de lichtere en bleekgroene tint van het
\VUgras{mos} \textsuperscript{(31)} waarin de \VUgras{groene specht}
\textsuperscript{(30)} insecten zoekt.

%%K t09-tit-11 sub
\VUsaut{3.0mm}
\VUpk{10.2pt}{40}{Woudlandschap.}
\VUsaut{3.0mm}

%%K t09-c2-07-1 p
7. --- De oude eiken bekoorden mij. Hun dikke gedraaide \VUgras{wortels}
\textsuperscript{(32)}, half uit de grond gekomen, geven hun een prachtig voorkomen van kracht
en sterkte, en ik vraag mij af hoe de \VUgras{eigenaar} \textsuperscript{(35)} zich ertoe kan
brengen ze te laten vellen en aan de
\VUgras{zaag} \textsuperscript{(34)} van de \VUgras{houthakker}
\textsuperscript{(33)} over te leveren. De arme \VUgras{stammen} \textsuperscript{(36)}, van
hun \VUgras{schors} \textsuperscript{(37)} beroofd of door de \VUgras{bijl}
\textsuperscript{(38)} van de \VUgras{dagloner} \textsuperscript{(39)} gespleten stemmen mij
droevig.

%%K t09-c2-08-1 p
8. --- Dichtbij had een oude \VUgras{houtskoolbrander} \textsuperscript{(41)} met zijn
\VUgras{schop een hoop} \textsuperscript{(42)} houtskool klaargemaakt, en een oude vrouw droeg
een \VUgras{bundel} \textsuperscript{(40)}
\VUgras{dood hout} van takjes van de gevelde bomen.

%%K t09-tit-12 sub
\VUsaut{3.0mm}
\VUpk{10.2pt}{40}{Jacht.}
\VUsaut{3.0mm}

%%K t09-c2-09-1 p
9. --- Ik wilde enkele ogenblikken gaan rusten in het \VUgras{huis van de boswachter}
\textsuperscript{(46)} maar, toen ik bij het \VUgras{meertje} \textsuperscript{(43)} kwam,
waarop een mooie \VUgras{zwaan} \textsuperscript{(45)} zwemt, merkte ik, dat een recente
\VUgras{overstroming} \textsuperscript{(44)} de weg in een waar \VUgras{moeras}
had veranderd. Ik moest een omweg maken, en, terwijl ik langs de \VUgras{ceder}
\textsuperscript{(49)}, de \VUgras{populieren} \textsuperscript{(48)}, en de
\VUgras{iep} \textsuperscript{(47)} kwam die aan de rand van het woud staan, zag
ik uit een \VUgras{kreupelbosje} \textsuperscript{(54)} een prachtig
\VUgras{hert} \textsuperscript{(50)} tevoorschijn komen, achtervolgd door een
hardnekkige \VUgras{meute} \textsuperscript{(51)} die ook aangespoord werd door de
\VUgras{hoorn} \textsuperscript{(53)} van de \VUgras{jagers te paard} \textsuperscript{(52)}.
Het arme dier was helemaal uitgeput. Het kwam stervend op enkele passen van mij neervallen.

%%K t09-c2-10-1 p
10. --- De zoon van de boswachter, die door het lawaai van de \VUgras{drijfjacht} aangetrokken
was, kwam het huis uit en wij gingen met ons tweeën het woud weer in. Hij had een
\VUgras{ekster} \textsuperscript{(56)} gezien op een tak van een oude eik. Hij meende, dat
haar nest waarschijnlijk in het \VUgras{loof} \textsuperscript{(58)} verborgen is en hij wilde
het nest pakken.

%%K t09-c2-10-2 p
Behendig als een \VUgras{eekhoorn} \textsuperscript{(61)}, klom hij van
\VUgras{tak} \textsuperscript{(55)} tot tak, maar hij vond niets en slaagde er
alleen in een oude \VUgras{kraai} \textsuperscript{(60)} te doen opvliegen die op een
\VUgras{dikke tak} \textsuperscript{(57)} zat.
\VUsaut{4.0mm}
\VUfilet{22mm}
